(1) 引力参数 \(\mu=GM>0\) 的点质量周围，一飞行器由半径 \(r_1\) 的顺行圆轨道转移至同平面半径 \(r_2>r_1\) 的顺行圆轨道。假设两次理想瞬时切向脉冲、两脉冲之间仅受牛顿万有引力，无扰动。求唯一在两圆轨道分别相切的椭圆转移轨道半长轴与两端速度。
(2) 求第一次加速脉冲 \(\Delta v_1\)、第二次加速脉冲 \(\Delta v_2\)，用角动量守恒独立核验近远地点速度，并说明轨道机械能在两次脉冲中的变化。
(3) 求半椭圆转移时间 \(t_t\)，与起始圆轨道周期比较，并检查 \(r_2\to r_1\) 和 \(r_2/r_1\to\infty\) 的物理极限。
(4) 若推进器有效喷气速度 \(v_e>0\) 恒定，两次脉冲之间无抛弃结构质量，求两次推进前后飞行器总质量比，并解释为什么总质量比不能由轨道机械能差直接代替速度增量求出。
